% Gauss's law on a sphere S' of radius R centred on a charge Q: E is parallel
% to dA everywhere and has the same size everywhere, which is what makes the
% flux integral trivial. The companion figure, ch01-gauss-surface-arbitrary.tex,
% shows the same charge inside an arbitrary surface S.
%
% Pulled in with \tikzfig{ch01-gauss-surface-sphere.tex}. A fragment, not a
% document -- one tikzpicture, no preamble -- so it inherits the fonts of
% whatever \inputs it. Styles match the companion figure, so the two read as a
% pair.
%
% Only E is drawn, as sampled on S' -- no field lines: equal arrows all round
% is the whole point of the sphere, and field lines would only repeat them.
%
% Layout: the eight E arrows sit at multiples of 45 deg, and everything that
% is lettered goes in the gaps between them, one thing per quadrant so the
% picture stays balanced --
%   upper right  E and dA, side by side at 45 deg (the one claim);
%   upper left   the name S', pointed at the sphere;
%   lower left   the radius R, inside the sphere;
%   right        the name Q, beside the charge.
% Each gap is named by its angle below, so a label can be moved to another
% quadrant by changing one number.
%
% Needs the arrows.meta and calc libraries, loaded in the preamble of
% gwbook.cls and _wrapper.tex -- see the note in ch01-coulomb-law.tex.
\begin{tikzpicture}[
    % The figsize knob: distances scale, lettering and node sizes do not.
    scale=1.0,
    >={Latex},
    charge/.style  = {circle, draw, fill=white, line width=0.9pt,
                      minimum size=12pt, inner sep=0pt, outer sep=2.5pt,
                      font=\normalsize},
    surface/.style = {line width=1.1pt},
    % See ch01-gauss-surface-arbitrary.tex.
    pointer/.style = {-{Straight Barb[round, length=4pt, width=4.6pt]},
                      line width=0.6pt, line cap=round, line join=round,
                      shorten >=3pt, shorten <=1pt},
    vec/.style     = {->, line width=0.9pt},
    % A measurement, not a field: lighter than the vectors, but with a head of
    % the same family so it does not look like an afterthought.
    % Not `radius': that would shadow the key circle[radius=...] needs, and
    % the sphere silently vanishes.
    radline/.style = {-{Latex[length=5pt, width=3.6pt]}, line width=0.7pt},
    % The patch of surface that dA stands for.
    patch/.style   = {line width=2.6pt, line cap=butt},
  ]

  \pgfmathsetmacro{\R}{1.7}       % the sphere
  \pgfmathsetmacro{\elen}{0.95}   % length of E
  \pgfmathsetmacro{\alen}{0.7}    % length of dA
  \pgfmathsetmacro{\half}{0.18}   % half-width of the dA patch
  \pgfmathsetmacro{\off}{2.8}     % pt; how far E and dA stand off their
                                  % shared normal, each to its own side
  % Where each lettered thing goes -- see the layout note above.
  \pgfmathsetmacro{\a}{45}        % E and dA
  \pgfmathsetmacro{\sname}{112.5} % S'
  \pgfmathsetmacro{\rdir}{202.5}  % R

  \node[charge] (Q) at (0,0) {$+$};
  \node[anchor=west, inner sep=2pt] at (Q.east) {$Q$};

  % E sampled on S': same length at every point, every one along the normal.
  % The one at \a is set aside to make room for dA beside it.
  \foreach \t in {0,45,...,315} {
    \ifdim \t pt=\a pt
      \draw[vec] ($(\t:\R)+(\t+90:\off pt)$) -- ++(\t:\elen)
        node[above left=-1pt] {$\mathbf{E}$};
    \else
      \draw[vec] (\t:\R) -- (\t:\R+\elen);
    \fi
  }
  \draw[vec] ($(\a:\R)+(\a-90:\off pt)$) -- ++(\a:\alen)
    node[below right=-1pt] {$d\mathbf{A}$};

  % The sphere's dashes are fitted to the circle -- 48 periods, 6 between
  % neighbouring arrows, phase starting mid-dash at 0 deg -- so every arrow
  % stands on a dash rather than in a gap. The period stays near the 6pt of
  % the dashes on S in the companion figure; keep the count a multiple of 8.
  % (dash pattern does not evaluate math, hence the macros.)
  \pgfmathsetmacro{\per}{2*pi*\R/48 * 28.4527559}   % cm -> pt
  \pgfmathsetmacro{\don}{\per*2/3}
  \pgfmathsetmacro{\doff}{\per/3}
  \pgfmathsetmacro{\dph}{\don/2}
  \draw[surface, dash pattern=on \don pt off \doff pt, dash phase=\dph pt]
    (0,0) circle[radius=\R];
  \draw[patch] ($(\a:\R)+(\a+90:\half)$) -- ($(\a:\R)+(\a-90:\half)$);

  % The radius, from Q's ring to S'. Lettered on the side away from Q's name.
  \draw[radline] (Q.\rdir) -- node[auto, inner sep=2pt] {$R$} (\rdir:\R);

  % S': the name sits out in its gap, a little round from the middle of it,
  % and the pointer swings back to land on the sphere mid-gap -- a gesture
  % rather than a radial tick, and clear of the arrows either side.
  \node (Slbl) at (\sname+10:\R+1.15) {$S'$};
  \draw[pointer] (Slbl.south) to[out=-85, in=160] (\sname-4:\R);

\end{tikzpicture}
